% !TEX root = ../../../main.tex
\section{Taylor's theorem}
\label{sec:analysis-i:taylor-theorem}

% BEGIN TUFTE PLACEHOLDER: supplied sample, lines 853--878.
\begingroup\sloppy
% Replace this entire specimen block with reviewed course notes.
\label{template:sec:typography}

\label{template:sec:typefaces}\index{typefaces}
If the Palatino, \textsf{Helvetica}, and \texttt{Bera Mono} typefaces are installed, this style
will use them automatically.  Otherwise, we'll fall back on the Computer Modern
typefaces.

\label{template:sec:letterspacing}
This document class includes two new commands and some improvements on
existing commands for letterspacing.

When setting strings of \allcaps{ALL CAPS} or \smallcaps{small caps}, the
letter\-spacing---that is, the spacing between the letters---should be
increased slightly.\cite[][p.32]{Bringhurst2005} The \doccmddef{allcaps} command has proper letterspacing for
strings of \allcaps{FULL CAPITAL LETTERS}, and the \doccmddef{smallcaps} command
has letterspacing for \smallcaps{small capital letters}.  These commands
will also automatically convert the case of the text to upper- or
lowercase, respectively.

The \doccmddef{textsc} command has also been redefined to include
letterspacing.  The case of the \doccmd{textsc} argument is left as is,
however.  This allows one to use both uppercase and lowercase letters:
\textsc{The Initial Letters Of The Words In This Sentence Are Capitalized.}

\lipsum[1]
\par\endgroup
% END TUFTE PLACEHOLDER
